
% Thesis Abstract -----------------------------------------------------


%\begin{abstractslong}    %uncommenting this line, gives a different abstract heading
\begin{abstracts}
\quad The development of sixth-generation (6G) networks aims to achieve global coverage through non-terrestrial networks (NTNs) utilizing low Earth orbit (LEO) satellites, with a focus on enhancing communication quality in high-traffic regions and areas lacking sufficient terrestrial infrastructure. However, the high altitude of LEO satellites increases the impact of signal fading, impairing accurate handover (HO) decisions by user equipment (UE). Additionally, the rapid mobility of the satellites results in frequent HOs, compromising network stability. This issue is particularly significant in Earth-moving cells (EMC) environments. Although conditional HO (CHO) and distance-based strategies have been introduced to address these challenges, they remain limited in effectively reducing unnecessary HOs (UHOs) and improving overall HO performance. To overcome these limitations, this thesis proposes an orbital information and distance (OID)-based CHO strategy. By using the positional relationship between the UE and satellite orbital information, the OID-based CHO enables more precise and efficient HO decisions, even in dynamic EMC environments where UE positions vary due to cell movement. Simulation results demonstrate that the OID-based CHO significantly minimizes UHOs and reduces signaling overhead compared to conventional approaches. Furthermore, the proposed strategy can be implemented without requiring additional information exchange, offering a practical and efficient solution for NTN-based LEO satellite communication in 6G NTN systems.

% The development of 6th generation (6G) networks focuses on non-terrestrial networks (NTNs) that provide worldwide coverage using low Earth orbit (LEO) satellites, especially in areas with high traffic or insufficient terrestrial infrastructure. However, frequent handovers (HO) occur due to the rapid mobility of LEO satellites and the effect of signal fading, which poses a problem for network stability, especially in Earth-moving cells (EMC) environments. Although distance-based strategies have been used to improve this, it is still difficult to improve HO performance, including unnecessary HOs (UHO). To resolve the challenge, this thesis proposes an orbital information and distance-based conditional HO (OID-based CHO) strategy that improves HO decision-making by utilizing the positional relationship between user equipment (UE) and satellite orbital information. This approach allows for more deliberate HO decision-making by considering variations in the UE's location status caused by cell movement in an EMC environment. Simulation results show that the OID-based CHO scheme minimizes UHOs and reduces signal overhead compared to conventional schemes. Furthermore, the proposed strategy does not require additional information exchange in NTN-based LEO satellite communication networks, providing a practical solution for 6G NTN systems.

% With the rise of 6th generation (6G) networks, Non-Terrestrial Networks (NTN) using Low Earth Orbit (LEO) satellites are crucial for global coverage, especially in underserved areas. However, managing handovers (HO) in LEO environments leads to frequent unnecessary handovers (UHO) and radio link failures (RLF). While 3GPP introduced Conditional HO (CHO) and distance-based strategies to address this, they can still result in suboptimal decisions in Earth-moving Cell (EMC) environments. This paper proposes an Orbital Information and Distance-based CHO (OID-based CHO) strategy, integrating real-time satellite orbital data with user equipment (UE) location. By leveraging satellite trajectory and position information, it minimizes UHO. Evaluating key metrics like downlink signal-to-interference-plus-noise ratio (DL SINR), UHO, RLF, and Time of Stay (ToS), results show that OID-based CHO reduces UHO, RLF, and signaling overhead, while enhancing ToS without large DL SINR reductions, thus providing better mobility management in LEO networks.

\end{abstracts} 

